\subsection{A bound below two for a prescribed five-site profile}
\label{sec:five-site-boundary}

Fixing the five selected values of the input $g$ gives a stronger
conclusion than the general five-site estimate. The restriction below
concerns these input values, with kernel $1_H$; it is not a restriction
on an arbitrary weighted kernel.

\begin{theorem}\label{thm:five-profile-boundary}
Let $H=\{h_0,h_1,h_2,h_3,h_4\}\subset\mathbb Z$ have five distinct
elements. Suppose $f,g$ are nonzero, nonnegative and finitely supported,
and, for some $t>0$,
\[
 g(h_0)=t,\qquad g(h_i)=\frac59t\quad(1\le i\le4).
\]
Then
\begin{equation}\label{eq:five-profile-boundary}
 \frac{\|f\star\widetilde{g1_H}\|_1}{\|f\star g\|_1}
 \le c_*:=\frac{27936}{13981}
 =2-\frac{26}{13981}<2.
\end{equation}
The positions $h_i$ and the values of $g$ outside $H$ are unrestricted.
\end{theorem}

The proof first treats all finite inputs on a free four-coordinate
lattice. A layer decomposition reduces arbitrary heights to weighted
sets. Counting their two boundaries, including overlaps on the reflected
side, gives the bound. Finite boxes in the kernel of the site map then
transfer it to arbitrary integer positions.

\begin{proof}
Put $r=5/9$, write $e_1,\ldots,e_4$ for the coordinate vectors in
$\mathbb Z^4$, and set $e_0=0$. For finite nonnegative $F$ on this
lattice, define
\[
\begin{aligned}
 P_0(F)&=\sum_x\max\{F(x),rF(x-e_1),\ldots,rF(x-e_4)\},\\
 Q_0(F)&=\sum_x\max\{F(x),rF(x+e_1),\ldots,rF(x+e_4)\}.
\end{aligned}
\]
We will prove $Q_0(F)<c_*P_0(F)$ for every nonzero such $F$.

\medskip\noindent\textbf{Reduction to weighted sets.}
Let $w(x)=r^{x_1+x_2+x_3+x_4}$, using the signed coordinate sum,
and set $u_S=w1_S$ for a finite set $S$. Decompose $F/w$ at its
distinct positive levels. This gives
\[
 F=\sum_{j=1}^m\lambda_j u_{S_j},\qquad \lambda_j>0,
\]
where the $S_j$ are nonempty finite nested sets. At every output $x$,
the five forward entries have the common factor $w(x)$, since
$rw(x-e_i)=w(x)$. An index maximizing the corresponding values of
$F/w$ also maximizes every threshold indicator. Consequently
\begin{equation}\label{eq:five-profile-layer-reduction}
 P_0(F)=\sum_j\lambda_jP_0(u_{S_j}),\qquad
 Q_0(F)\le\sum_j\lambda_jQ_0(u_{S_j}).
\end{equation}
The second assertion is subadditivity of the maximum. It therefore
suffices to prove the strict bound for every nonempty finite $S$.
No monotonicity or shape condition is imposed on these sets.

\medskip\noindent\textbf{The two boundaries.}
For finite $A$, write $\mu(A)=\sum_{x\in A}w(x)$ and put $m=\mu(S)$.
For $x\notin S$, let
\[
 k_-(x)=\#\{i:x-e_i\in S\},\qquad
 k_+(x)=\#\{i:x+e_i\in S\},
\]
where $1\le i\le4$. For $1\le j\le4$ define the finite boundary masses
\[
 U_j=\mu\{x\notin S:k_-(x)=j\},\qquad
 V_j=\mu\{x\notin S:k_+(x)=j\}.
\]
All nonzero forward entries at $x$ equal $w(x)$. A reflected leaf
equals $r^2w(x)$, and the root dominates on $S$. Thus
\[
 P_0(u_S)=m+\sum_jU_j,\qquad Q_0(u_S)=m+r^2\sum_jV_j.
\]
Set $I_i=\sum_{x\in S}w(x)1_S(x-e_i)$. In direction $i$, the
forward outside mass is $rm-I_i$, and the reflected outside mass
is $(m-I_i)/r$. Summing over the four directions gives
\[
 r\sum_jjV_j-\sum_jjU_j=4(1-r)m.
\]
The reflected overlap mass $E=\sum_j(j-1)V_j$ is nonnegative, and
the preceding identities give
\begin{equation}\label{eq:five-profile-boundaries}
 2P_0(u_S)-Q_0(u_S)
 =\frac{m}{81}+\frac{13}{9}U_1+\frac89U_2+\frac13U_3
       -\frac29U_4+\frac{25}{81}E.
\end{equation}
In particular, a reflected boundary point shared by several translates
contributes its maximum once; $E$ accounts for the excess multiplicity.

\medskip\noindent\textbf{Complete lower cubes.}
For $J\subseteq\{1,2,3,4\}$, put $1_J=\sum_{i\in J}e_i$.
Partition the points $y\notin S$ with $k_-(y)=4$ into two classes.
Call $y$ complete if $y-1_J\in S$ for every nonempty $J$; otherwise
call it incomplete. Let $G$ and $B$ be their respective weighted
masses, so $U_4=G+B$.

For each $x\in S$, the nonempty subsets $J$ for which $x+1_J$ is
complete form an antichain. If $J\subsetneq K$ both qualified,
completeness of $x+1_K$ would put
$(x+1_K)-1_{K\setminus J}=x+1_J$ in $S$, a contradiction.
A random permutation of the four directions has a given size-$j$
subset as its first $j$ directions with probability $\binom4j^{-1}$.
Its nested prefix sets meet an antichain at most once, proving
\[
 \sum_{\substack{J\ne\varnothing\\x+1_J\text{ complete}}}
          \binom4{|J|}^{-1}\le1.
\]
Multiply by $w(x)$ and sum over $x\in S$. Each complete point $y$
contributes from all its nonempty lower corners. The binomial factors
cancel the number of corners of each size, yielding
\begin{equation}\label{eq:five-profile-complete-cubes}
 m\ge aG,\qquad a=\sum_{j=1}^4r^{-j}=\frac{13356}{625}.
\end{equation}

\medskip\noindent\textbf{Incomplete lower cubes.}
For each incomplete point $y$, choose a nonempty $J$ of minimum
cardinality with $z=y-1_J\notin S$, breaking ties by a fixed rule.
Since all four immediate predecessors of $y$ belong to $S$,
$j=|J|\ge2$. Minimality gives $z+e_i\in S$ for every $i\in J$.
Thus $J$ is a subset of the $k_+(z)$ successor directions at $z$,
and $w(y)=r^jw(z)$.

At a fixed $z$ with $k=k_+(z)$, distinct assigned points $y$ have
distinct subsets $J$, since $y=z+1_J$. Their total weight is at most
$w(z)\sum_{j=2}^k\binom kjr^j$. The nonzero coefficients are
\[
\begin{array}{c|ccc}
 k&2&3&4\\ \hline
 \displaystyle\sum_{j=2}^k\binom kjr^j
       &25/81&800/729&17275/6561
\end{array}
\]
and each is at most $k-1$. There is no assignment for $k=0,1$.
Summing over missing points proves
\begin{equation}\label{eq:five-profile-incomplete-cubes}
 B\le\sum_{z\notin S}(k_+(z)-1)_+w(z)=E.
\end{equation}
Several incomplete points may have the same assigned missing point;
the binomial sum includes all such assignments.

\medskip\noindent\textbf{Combining the counts.}
Put $\varepsilon=26/13981$. Then $0<\varepsilon<1/81$ and
$a(1/81-\varepsilon)=2/9+\varepsilon$. Substitute $U_4=G+B$
in \eqref{eq:five-profile-boundaries} to obtain the exact identity
\[
\begin{aligned}
 (2-\varepsilon)P_0(u_S)-Q_0(u_S)
 ={}&(1/81-\varepsilon)(m-aG)\\
    &+(13/9-\varepsilon)U_1+(8/9-\varepsilon)U_2
       +(1/3-\varepsilon)U_3\\
    &+(2/9+\varepsilon)(E-B)+(7/81-\varepsilon)E.
\end{aligned}
\]
Every coefficient is positive and every multiplied quantity is
nonnegative, by \eqref{eq:five-profile-complete-cubes} and
\eqref{eq:five-profile-incomplete-cubes}. Moreover $U_1>0$:
choose $x\in S$ with maximal first coordinate. The point $x+e_1$
lies outside $S$ and has exactly one predecessor in $S$, because
all its other possible predecessors have a larger first coordinate.
This proves the strict set bound and, by
\eqref{eq:five-profile-layer-reduction},
\begin{equation}\label{eq:five-profile-free-bound}
 Q_0(F)<c_*P_0(F)\qquad(0\ne F\ge0\text{ finitely supported}).
\end{equation}

\medskip\noindent\textbf{Transfer to integer positions.}
Translate the selected sites by $-h_0$, writing $s_i=h_i-h_0$.
Put $s_0=0$, $a_0=1$, and $a_i=r$ for $1\le i\le4$. For any
integer shift list $s$ define
\[
 P_s(u)=\sum_n\max_{0\le i\le4}a_iu(n-s_i),\qquad
 Q_s(u)=\sum_n\max_{0\le i\le4}a_iu(n+s_i).
\]
Let $b=\gcd(|s_1|,\ldots,|s_4|)>0$ and $d_i=s_i/b$.
Splitting $f$ into its residue classes modulo $b$ reduces the desired
inequality to $Q_d(u)\le c_*P_d(u)$ for each finite nonnegative
residue function $u$.

The map $\varphi:\mathbb Z^4\to\mathbb Z$ given by
$\varphi(x)=\sum_i d_ix_i$ is surjective. Bezout's identity gives
$v\in\mathbb Z^4$ with $\varphi(v)=1$, so
$\mathbb Z^4=\mathbb Z v\oplus\ker\varphi$.
The kernel has an integer basis $b_1,b_2,b_3$; for example, apply
unimodular Euclidean operations to the primitive row $(d_i)$.
Write $Bm=\sum_{j=1}^3m_jb_j$ for $m\in\mathbb Z^3$, and choose
$\alpha_i\in\mathbb Z^3$ with $e_i=d_iv+B\alpha_i$; set
$d_0=0$ and $\alpha_0=0$.

For $C_N=\{m\in\mathbb Z^3:|m_j|\le N\}$, define the finite lift
\[
 F_N(nv+Bm)=u(n)1_{C_N}(m).
\]
Let $M=\max_{i,j}|(\alpha_i)_j|$. At $x=nv+Bm$ the two maxima
in the free-lattice sums are
\[
 \max_i a_iu(n-d_i)1_{C_N}(m-\alpha_i),\qquad
 \max_i a_iu(n+d_i)1_{C_N}(m+\alpha_i).
\]
For $N\ge M$, these equal their untruncated one-dimensional maxima
throughout $C_{N-M}$, vanish outside $C_{N+M}$, and never exceed
those maxima elsewhere. Therefore
\[
\begin{aligned}
 |C_{N-M}|P_d(u)&\le P_0(F_N)\le|C_{N+M}|P_d(u),\\
 |C_{N-M}|Q_d(u)&\le Q_0(F_N)\le|C_{N+M}|Q_d(u).
\end{aligned}
\]
Since $|C_N|=(2N+1)^3$, division by $|C_N|$ and passage to the
limit in \eqref{eq:five-profile-free-bound} give
$Q_d(u)\le c_*P_d(u)$. Zero residue functions contribute zero,
and summing the residue inequalities gives $Q_s(f)\le c_*P_s(f)$.
This argument permits overlaps and relations among the integer shifts.
The strict free-lattice inequality yields a non-strict limiting bound.

Finally, translating the summation variable separately in the forward
and reflected totals gives
\[
 \|f\star(g1_H)\|_1=tP_s(f),\qquad
 \|f\star\widetilde{g1_H}\|_1=tQ_s(f).
\]
The full denominator satisfies
$\|f\star g\|_1\ge tP_s(f)>0$, proving
\eqref{eq:five-profile-boundary}.
\end{proof}

This bound does not determine the best constant for the prescribed
profile, or the general five-site constant $K_5$. It concerns the
reflected ratio, not every intermediate estimate used to bound that
ratio. The following certificate calculation retains separate
information about those estimates and their unused margins.
